\section{Introducción}

El maíz (\textit{Zea mays}) constituye uno de los cultivos de mayor relevancia a nivel global, tanto desde el punto de vista productivo como científico. En este contexto, resulta de interés explorar intervenciones físicas de bajo costo que puedan influir sobre la germinación y el crecimiento inicial de las semillas.

Diversos trabajos \cite{araujo_2016,sarraf_2020,vashisth_joshi_2017,ferroni_2024} han reportado que los campos magnéticos estáticos pueden actuar como moduladores de procesos biológicos en distintas especies vegetales. Sin embargo, los resultados obtenidos en la literatura muestran una alta dependencia de las condiciones experimentales, tales como la intensidad del campo, su geometría, la orientación relativa y los tiempos de exposición. Esta variabilidad dificulta la identificación de efectos robustos y reproducibles.

En este trabajo se propone abordar este problema mediante un diseño experimental controlado, orientado a la exploración sistemática de parámetros de exposición. En particular, se busca responder a la siguiente pregunta de investigación:

\begin{quote}
¿Existe un rango de condiciones de exposición a campos magnéticos estáticos que produzca mejoras reproducibles en la germinación y el crecimiento inicial de \textit{Zea mays} respecto de un grupo control?
\end{quote}

La hipótesis de trabajo plantea que, bajo ciertas combinaciones de parámetros físicos de exposición, se observarán cambios medibles y reproducibles en indicadores de germinación y crecimiento temprano.

El objetivo general consiste en caracterizar el impacto de condiciones físicas de exposición, incluyendo campos magnéticos estáticos, sobre la germinación y el crecimiento inicial de semillas de \textit{Zea mays}, identificando tendencias, rangos efectivos y su reproducibilidad bajo un protocolo experimental controlado.

Para ello, se priorizará la definición de un protocolo reproducible, la validación de un sistema de medición basado en imágenes (photobox y procesamiento algorítmico), y la cuantificación rigurosa de los efectos observados mediante herramientas estadísticas.

\clearpage
